\documentclass[12pt]{article}
\usepackage[utf8]{inputenc}
\usepackage[T1]{fontenc}
\usepackage{hyperref}
\title{Nunukul Language Resources for K-12 Curriculum Development}
\author{Compiled by Claude Sonnet 5.5 (claude-sonnet-5-5) for Andr\'as Kornai}
\date{}
\begin{document}
\maketitle

\section{Introduction}
Nunukul (ISO 639-3 xnu) is an extinct Australian language. The Nunukul (also Noonuccal or Moondjan) are one of three Quandamooka groups, with the Goenpul and Ngugi, and lived on Minjerribah (North Stradbroke Island) and in the Moreton Bay region. Sources disagree on the language's name and classification. Some treat it as a dialect of Yagara or Durubalic, called Moondjan or Mundjan, and AIATSIS lists it as Moondjan (E21). The search found very little material specific to Nunukul.

\section{Primary Educational Resources}
The search found no primers, textbooks or school programmes for Nunukul.

\section{Linguistic and Reference Materials}
\begin{itemize}
\item Nunukul, Wikipedia: \url{https://en.wikipedia.org/wiki/Nunukul}
\item Quandamooka people, Wikipedia: \url{https://en.wikipedia.org/wiki/Quandamooka_people}
\item AIATSIS Austlang record E21 (Moondjan): \url{https://collection.aiatsis.gov.au/austlang/language/e21}
\item AIATSIS collection record: \url{https://collection.aiatsis.gov.au/node/1508184}
\item Janday language, Wikipedia. It is a related Moreton Bay language, with a note on Dixon's grouping of Yagara dialects: \url{https://Www.Wikipedia.org/wiki/Jandai}
\end{itemize}

\section{Religious and Community Resources}
The search found no religious resources. Several web pages titled ``Quandamooka Country Language'' appeared, but they looked garbled or machine-generated and are not cited.

\section{Digital and Audio Resources}
No audio, video or software resources were found.

\section{Curriculum Development Context}
Little is documented for Nunukul. Any heritage curriculum would need to work with Quandamooka community organisations and with AIATSIS archival holdings, and might draw on better-documented neighbouring Yagara and Janday materials. Which of those materials apply to Nunukul would need expert advice. The identification of Nunukul with Moondjan or Yagara should be checked with the community.

\end{document}